% !TEX program = latexmk
% Compile example:
% latexmk -pdfdvi -latex='uplatex -synctex=1 -interaction=nonstopmode -file-line-error' -dvipdf='dvipdfmx %O -o %D %S' d2l_why_conv_seminar_beamer.tex

\documentclass[dvipdfmx,uplatex,unicode,9pt,aspectratio=169]{beamer}

% ------------------------------------------------------------
% Japanese / PDF settings
% ------------------------------------------------------------
\usepackage{bxdpx-beamer}
\usepackage{pxjahyper}
\usepackage{amsmath,amssymb,bm,mathtools}
\usepackage{tikz}
\usetikzlibrary{positioning,arrows.meta,calc,decorations.pathreplacing,fit,matrix}
\usepackage{booktabs}
\usepackage{url}
\usepackage{listings}
\lstdefinestyle{pytorch}{
  language=Python,
  basicstyle=\ttfamily\scriptsize,
  keywordstyle=\color{blue!70!black}\bfseries,
  commentstyle=\color{green!40!black},
  stringstyle=\color{orange!60!black},
  showstringspaces=false,
  breaklines=true,
  frame=single,
  framerule=0.2pt,
  rulecolor=\color{black!25},
  backgroundcolor=\color{black!3},
  columns=fullflexible,
  keepspaces=true,
  xleftmargin=2mm,
  xrightmargin=2mm
}
\lstset{style=pytorch}

% ------------------------------------------------------------
% Custom boxes
% ------------------------------------------------------------
\usepackage[most]{tcolorbox}

\tcbset{
  seminarboxstyle/.style={
    enhanced,
    colback=softblue,
    colframe=mainblue,
    boxrule=0.8pt,
    arc=2mm,
    left=1.5mm,
    right=1.5mm,
    top=1.2mm,
    bottom=1.2mm,
    before skip=0.6em,
    after skip=0.6em,
    fonttitle=\bfseries,
    coltitle=white,
    colbacktitle=mainblue,
    attach boxed title to top left={
      xshift=2mm,
      yshift=-2mm
    },
    boxed title style={
      colback=mainblue,
      colframe=mainblue,
      arc=1mm,
      boxrule=0pt
    }
  }
}

% タイトルあり・なしの両方で使える箱
\newtcolorbox{seminarbox}[1][]{
  seminarboxstyle,
  #1
}

% 注意用のオレンジ箱
\newtcolorbox{seminaralertbox}[1][]{
  enhanced,
  colback=softorange,
  colframe=accentorange,
  boxrule=0.8pt,
  arc=2mm,
  left=1.5mm,
  right=1.5mm,
  top=1.2mm,
  bottom=1.2mm,
  before skip=0.6em,
  after skip=0.6em,
  fonttitle=\bfseries,
  coltitle=white,
  colbacktitle=accentorange,
  attach boxed title to top left={
    xshift=2mm,
    yshift=-2mm
  },
  boxed title style={
    colback=accentorange,
    colframe=accentorange,
    arc=1mm,
    boxrule=0pt
  },
  #1
}

% ------------------------------------------------------------
% Beamer style
% ------------------------------------------------------------
\usetheme{Madrid}
\usefonttheme{professionalfonts}
\setbeamertemplate{navigation symbols}{}
\setbeamertemplate{footline}[frame number]

\definecolor{mainblue}{RGB}{36,84,150}
\definecolor{accentorange}{RGB}{226,122,47}
\definecolor{softblue}{RGB}{232,240,252}
\definecolor{softorange}{RGB}{255,239,224}
\definecolor{softgray}{RGB}{245,245,245}
\setbeamercolor{structure}{fg=mainblue}
\setbeamercolor{frametitle}{fg=white,bg=mainblue}
\setbeamercolor{title}{fg=white,bg=mainblue}
\setbeamercolor{block title}{fg=white,bg=mainblue}
\setbeamercolor{block body}{bg=softblue}
\setbeamercolor{alerted text}{fg=accentorange}

\newcommand{\R}{\mathbb{R}}
\newcommand{\X}{\mathbf{X}}
\newcommand{\Y}{\mathbf{Y}}
\newcommand{\K}{\mathbf{K}}
\newcommand{\Hid}{\mathbf{H}}
\newcommand{\V}{\mathbf{V}}
\newcommand{\Wten}{\mathsf{W}}
\newcommand{\Vten}{\mathsf{V}}
\newcommand{\Xten}{\mathsf{X}}
\newcommand{\Hten}{\mathsf{H}}
\newcommand{\U}{\mathbf{U}}
\newcommand{\patch}{\mathcal{P}}

\title[全結合層から畳み込みへ]{CNNとは：全結合層から畳み込みへ}
\subtitle{Dive into Deep Learning 輪読ゼミ}
\author{山口 海}
\date{\today}

\begin{document}

% ------------------------------------------------------------
\begin{frame}
  \titlepage
  \vspace{-1.5em}
  \begin{center}
    \footnotesize 参考: D2L-ja 7.1, \url{https://d2l-jp.me/chapter_convolutional-neural-networks/why-conv.html}
  \end{center}
\end{frame}

%================================================
\begin{frame}[allowframebreaks]{目次}
  \tableofcontents[
    sectionstyle=show,
    subsectionstyle=show/show/show
  ]
\end{frame}

%================================================
\section{7.1. CNNとは: 全結合層から畳み込みへ}
\begin{frame}{7.1. CNNとは: 全結合層から畳み込みへ}
\begin{block}{問い}
なぜ画像認識では，全結合層ではなく\alert{畳み込み層}を使うのか？
\end{block}
\begin{columns}[T,onlytextwidth]
\begin{column}{0.48\textwidth}
\begin{enumerate}
 \item 画像は高次元で，そのまま全結合にするとパラメータが爆発する．
  \item しかし画像には，位置関係や局所構造という強い規則性がある．
  \item その規則性を\alert{帰納バイアス}としてモデルに入れる．
  \item その結果として，自然に\alert{畳み込み}が現れる．
\end{enumerate}
\end{column}
\begin{column}{0.44\textwidth}
\begin{center}
\begin{tikzpicture}[scale=0.58]
\foreach \x in {0,...,4}{\foreach \y in {0,...,4}{\fill[mainblue!25] (\x,\y) rectangle +(0.85,0.85);}}
\node at (2.2,-0.45) {画像：多数の画素};
\foreach \y in {0,...,3}{\fill[accentorange!40] (6,\y+0.4) circle (0.18);}
\node at (6,-0.45) {隠れ層};
\foreach \x in {0,...,4}{\foreach \y in {0,...,4}{\draw[mainblue!25] (\x+0.42,\y+0.42) -- (6,2.0);}}
\node[align=center] at (3.8,5.0) {全部つなぐと\\重みが爆発};
\end{tikzpicture}
\end{center}
\end{column}
\end{columns}

\begin{center}
\begin{tikzpicture}[node distance=13mm, >=Latex]
\node[draw, rounded corners, fill=softgray, minimum width=26mm, minimum height=9mm] (mlp) {全結合MLP};
\node[draw, rounded corners, fill=softorange, right=of mlp, minimum width=26mm, minimum height=9mm] (bias) {画像の構造};
\node[draw, rounded corners, fill=softblue, right=of bias, minimum width=28mm, minimum height=9mm] (conv) {畳み込み層};
\draw[->, thick] (mlp) -- node[above]{制約} (bias);
\draw[->, thick] (bias) -- node[above]{導出} (conv);
\end{tikzpicture}
\end{center}
\end{frame}

\begin{frame}{畳み込み層の計算イメージ}
\begin{center}
\begin{tikzpicture}[scale=0.63, >=Latex]
% input grid
\foreach \x in {0,...,5}{\foreach \y in {0,...,5}{\draw[mainblue!35] (\x,\y) rectangle +(0.8,0.8);}}
\foreach \x in {2,3,4}{\foreach \y in {2,3,4}{\fill[accentorange!30] (\x,\y) rectangle +(0.8,0.8);\draw[accentorange] (\x,\y) rectangle +(0.8,0.8);}}
\node at (2.4,-0.7) {入力画像 $\X$};
\node[accentorange] at (3.2,5.4) {局所パッチ};

% kernel
\foreach \x in {0,...,2}{\foreach \y in {0,...,2}{\draw[accentorange] (7+\x,2+\y) rectangle +(0.8,0.8);}}
\node at (8.2,1.3) {カーネル $\V$};

% output grid
\foreach \x in {0,...,3}{\foreach \y in {0,...,3}{\draw[mainblue!35] (12+\x,1+\y) rectangle +(0.8,0.8);}}
\fill[mainblue!45] (13,2) rectangle +(0.8,0.8);
\node at (13.6,0.3) {特徴マップ $\Hid$};

\draw[->, thick] (5.3,3.2) -- node[above]{重み付き和} (6.8,3.2);
\draw[->, thick] (9.8,3.2) -- node[above]{1つの値} (11.8,2.5);
\end{tikzpicture}
\end{center}

\vspace{2mm}
\begin{block}{1つの出力値の意味}
入力画像の小さなパッチとカーネルの「一致度」を測ったスコアと考えられる．
\end{block}
\end{frame}

%------------------------------------------------------------
\begin{frame}{CNNに入れる2つの帰納バイアス}
\begin{columns}[T,onlytextwidth]
\begin{column}{0.48\textwidth}
\begin{block}{1. 平行移動等価性 / 不変性}
同じ局所パターンが画像のどこに現れても，同じように反応してほしい．
\[
\text{位置がずれる} \Rightarrow \text{特徴マップもずれる}
\]
\end{block}
\end{column}
\begin{column}{0.48\textwidth}
\begin{block}{2. 局所性}
ある位置の特徴を計算するには，画像全体ではなく，その近くの画素だけ見ればよい．
\[
\text{近傍パッチ} \Rightarrow \text{局所特徴}
\]
\end{block}
\end{column}
\end{columns}

\vspace{4mm}
\begin{center}
\begin{tikzpicture}[>=Latex, scale=0.9]
\node[draw, rounded corners, fill=softorange, minimum width=30mm, minimum height=9mm] (eq) at (0,0) {同じ検出器を使い回す};
\node[draw, rounded corners, fill=softblue, minimum width=28mm, minimum height=9mm] (loc) at (5.2,0) {近所だけを見る};
\node[draw, rounded corners, fill=softgray, minimum width=34mm, minimum height=9mm] (conv) at (2.6,-1.7) {畳み込み層};
\draw[->, thick] (eq) -- (conv);
\draw[->, thick] (loc) -- (conv);
\end{tikzpicture}
\end{center}
\begin{block}{この節で理解したいこと}
画像にMLPをそのまま使うとパラメータが多すぎる. しかし画像には,
\alert{平行移動に対する性質}と\alert{局所性}がある. これらをMLPに制約として入れると,
自然に\alert{畳み込み層}が出てくる.
\end{block}

\end{frame}

% ------------------------------------------------------------
\subsection{7.1.1. 不変性}

\begin{frame}{7.1.1. 不変性: どこにあっても同じものは同じ}
\begin{seminarbox}
猫, 縦線, 目のような特徴は, 画像の左上にあっても右下にあっても,
同じ種類の特徴として認識したい.
\end{seminarbox}

\begin{columns}[T,totalwidth=\textwidth]
\begin{column}{0.48\textwidth}
\centering
\begin{tikzpicture}[scale=0.55]
  \draw[step=0.8,gray!60] (0,0) grid (5.6,4.0);
  \draw[fill=accentorange!70] (0.8,2.4) rectangle (1.6,3.2);
  \node at (2.8,-0.5) {左上に特徴};
\end{tikzpicture}
\end{column}
\begin{column}{0.48\textwidth}

\begin{tikzpicture}[scale=0.55]
  \draw[step=0.8,gray!60] (0,0) grid (5.6,4.0);
  \draw[fill=accentorange!70] (4.0,0.8) rectangle (4.8,1.6);
  \node at (2.8,-0.5) {右下に特徴};
\end{tikzpicture}
\end{column}
\end{columns}
\begin{itemize}
  \item 最終的な分類では, 多少位置が変わっても同じクラスと判断したい.
  \item 中間層では, 入力をずらすと特徴マップの反応もずれる, という性質が欲しい.
\end{itemize}
\begin{columns}[T,totalwidth=\textwidth]
\begin{column}{0.48\textwidth}
\begin{block}{平行移動不変性}
入力をずらしても, 出力が変わらない.
\[
  f(T\X)=f(\X)
\]
\end{block}
\begin{itemize}
  \item $\text{猫画像} \longmapsto \text{猫}$
  \item $\text{右にずれた猫画像} \longmapsto \text{猫}$
\end{itemize}
\end{column}
\begin{column}{0.48\textwidth}
\begin{block}{平行移動等価性}
入力をずらすと, 出力も同じようにずれる.
\[
  f(T\X)=T f(\X)
\]
\end{block}
\begin{itemize}
  \item 特徴が移動すると, 特徴マップ上の反応位置も移動する.
\end{itemize}
\end{column}
\end{columns}
\end{frame}

% ------------------------------------------------------------


% ------------------------------------------------------------


% ------------------------------------------------------------
\subsection{7.1.2. MLPへの制約}
\begin{frame}{7.1.2. MLPへの制約: 画像を2次元配列として扱う}
\begin{columns}[T,onlytextwidth]
\begin{column}{0.63\textwidth}
\begin{block}{ベクトル入力の全結合層}
入力を $\bm{x}=(x_1,\dots,x_d)^\top$, 隠れ表現を $\bm{h}=(h_1,\dots,h_m)^\top$ とすると,
\[
  \bm{h}=\bm{W}\bm{x}+\bm{b}
\]
\[
  h_i=b_i+\sum_{k=1}^{d}W_{i,k}x_k.
\]
\end{block}
\end{column}
\begin{column}{0.37\textwidth}
\begin{itemize}
  \item $h_i$ は隠れ層の $i$ 番目のユニット.
  \item $x_k$ は入力の $k$ 番目の成分.
  \item $W_{i,k}$ は「入力 $k$ が隠れユニット $i$ に与える重み」.
\end{itemize}
\begin{seminaralertbox}
全結合層では, 各出力 $[\Hid]_{i,j}$  がすべての入力画素 $[\X]_{k,l}$ を見る.
\end{seminaralertbox}
\end{column}
\end{columns}

\vspace{0.6em}

\begin{columns}[T,totalwidth=\textwidth]
\begin{column}{0.52\textwidth}
\begin{block}{ベクトルから行列へ}
出力の位置が $(i,j)$, 入力の位置が $(k,l)$ なので,
\[
  [\Hid]_{i,j}
  = [\U]_{i,j}
  +\sum_k\sum_l [\Wten]_{i,j,k,l}[\X]_{k,l}.
\]
\end{block}

\vspace{-0.3em}

\begin{itemize}
  \item ここでは簡単のため, $\X$ と $\Hid$ は同じ形状とする.
  \item $[\U]_{i,j}$ は位置 $(i,j)$ のバイアス.
\end{itemize}
\end{column}
\begin{column}{0.44\textwidth}
\centering
\begin{tikzpicture}[scale=0.55]
  % 入力画像 X
  \draw[step=0.75,gray!65] (0,0) grid (4.5,3.75);
  \node at (2.25,4.2) {$\X$};
  \node[fill=accentorange!60,inner sep=2pt] at (2.625,2.625)
    {$[\X]_{k,l}$};

  % X から H への写像
  \draw[-{Latex[length=2mm]},thick]
    (5.0,1.875) -- (6.3,1.875)
    node[midway,above=3pt] {$\Wten$};

  % 隠れ表現 H
  \draw[step=0.75,gray!65] (6.8,0) grid (11.3,3.75);
  \draw[gray!65] (6.8,0) rectangle (11.3,3.75);
  \node at (9.05,4.2) {$\Hid$};
  \node[fill=softblue,inner sep=2pt] at (9.425,2.625)
    {$[\Hid]_{i,j}$};
\end{tikzpicture}

入力画像と隠れ表現を，2次元の行列として扱う．
\[
  \X = \bigl([\X]_{i,j}\bigr),\qquad
  \Hid = \bigl([\Hid]_{i,j}\bigr)
\]
\end{column}
\end{columns}
\end{frame}



% ------------------------------------------------------------


% ------------------------------------------------------------
\begin{frame}{添字を相対位置で書き直す}
\begin{columns}[T,totalwidth=\textwidth]
\begin{column}{0.55\textwidth}
\begin{block}{パラメータ数の爆発}
全結合では,
\[
  \text{入力画素数}\times \text{出力ユニット数}
  =10^6\times10^6=10^{12}.
\]
\end{block}
\end{column}
\begin{column}{0.41\textwidth}
\centering
\begin{tikzpicture}[scale=0.75]
  \foreach \i in {0,...,3}{
    \foreach \j in {0,...,3}{
      \node[circle,draw,fill=softblue,minimum size=5pt,inner sep=0pt] (x\i\j) at (\i*0.6,\j*0.55) {};
      \node[circle,draw,fill=softorange,minimum size=5pt,inner sep=0pt] (h\i\j) at (3.6+\i*0.6,\j*0.55) {};
    }
  }
  \foreach \i in {0,...,3}{
    \foreach \j in {0,...,3}{
      \foreach \k in {0,...,3}{
        \foreach \l in {0,...,3}{
          \draw[gray!25] (x\i\j) -- (h\k\l);
        }
      }
    }
  }
  \node at (0.9,2.55) {入力 $\X$};
  \node at (4.5,2.55) {隠れ表現 $\Hid$};
  \node[accentorange] at (2.7,-0.55) {全結合は接続が多すぎる};
\end{tikzpicture}
\end{column}
\end{columns}
\begin{block}{入力位置を出力位置からのズレで表す}
出力位置 $(i,j)$ に対して, 入力位置 $(k,l)$ を
\[
  k=i+a,\qquad l=j+b
\]
と書く. ここで $(a,b)$ は出力位置から見た相対的なズレである.
\end{block}
\vspace{-0.4em}
\begin{columns}[T,totalwidth=\textwidth]
\begin{column}{0.55\textwidth}
  \centering
\begin{align*}
  [\Hid]_{i,j}
  &= [\U]_{i,j}
  +\sum_k\sum_l [\Wten]_{i,j,k,l}[\X]_{k,l} \\
  &= [\U]_{i,j}
  +\sum_a\sum_b [\Vten]_{i,j,a,b}[\X]_{i+a,j+b}.
\end{align*}
\end{column}
\begin{column}{0.41\textwidth}
\begin{block}{対応関係}
\[
  [\Vten]_{i,j,a,b}
  = [\Wten]_{i,j,i+a,j+b}.
\]
\end{block}
\end{column}
\end{columns}
\end{frame}

% ------------------------------------------------------------
\subsection{7.1.2.1. 平行移動不変性}
\begin{frame}{7.1.2.1. 平行移動不変性}
\begin{block}{重み共有: 平行移動に対する制約}
重みとバイアスが出力位置 $(i,j)$ に依存しないと仮定する.
\[
  [\Vten]_{i,j,a,b}=[\V]_{a,b},
  \qquad
  [\U]_{i,j}=u.
\]
すると,
\[
  [\Hid]_{i,j}
  =u+\sum_a\sum_b [\V]_{a,b}[\X]_{i+a,j+b}.
\]
\end{block}
\begin{columns}[T,totalwidth=\textwidth]
\begin{column}{0.52\textwidth}
\begin{itemize}
  \item $[\V]_{a,b}$は相対位置 $(a,b)$ だけで重みが決まる.
  \item 同じ特徴検出器を画像全体に適用する.
  \item これを\alert{重み共有}と呼ぶ.
\end{itemize}
\end{column}
\begin{column}{0.42\textwidth}
\centering
\begin{tikzpicture}[scale=0.6]
  \draw[step=0.75,gray!60] (0,0) grid (5.25,3.75);
  \foreach \x/\y in {0.75/2.25,2.25/1.5,3.75/0.75}{
    \draw[very thick,accentorange] (\x,\y) rectangle +(1.5,1.5);
  }
  \node at (2.6,-0.5) {同じフィルタを使い回す};
\end{tikzpicture}
\end{column}
\end{columns}
\begin{block}{まだ残る問題}
それでも, 各位置で画像全体に近い広い範囲を見ている. そこで次に局所性を入れる.
\end{block}
\end{frame}
% ------------------------------------------------------------
\subsection{7.1.2.2. 局所性}
\begin{frame}{7.1.2.2. 局所性を式に入れる}
\begin{block}{局所範囲だけを残す}
ある $\Delta$ を決めて,
\[
  |a|>\Delta \quad \text{または}\quad |b|>\Delta
  \quad \Longrightarrow \quad
  [\V]_{a,b}=0
\]
とする.
\end{block}

すると,
\[
  [\Hid]_{i,j}
  =u+
  \sum_{a=-\Delta}^{\Delta}
  \sum_{b=-\Delta}^{\Delta}
  [\V]_{a,b}[\X]_{i+a,j+b}.
\]
\begin{columns}[T,totalwidth=\textwidth]
\begin{column}{0.5\textwidth}
\centering
\begin{tikzpicture}[scale=0.58]
  \draw[step=0.7,gray!60] (0,0) grid (6.3,4.2);
  \draw[fill=accentorange!15,very thick,accentorange] (2.1,1.4) rectangle (4.2,3.5);
  \node[fill=mainblue!80,text=white,inner sep=1.5pt] at (3.15,2.45) {$(i,j)$};
  \node[accentorange] at (3.15,4.65) {近傍だけを見る};
\end{tikzpicture}
\end{column}
\begin{column}{0.46\textwidth}
\begin{itemize}
  \item 相対位置 $(a,b)$ が大きすぎるとき, 重みを0にする.
  \item エッジや模様は局所パッチから分かる.
  \item 遠く離れた画素との相互作用は, 深い層で段階的に扱えばよい.
\end{itemize}
\end{column}
\end{columns}
\end{frame}

% ------------------------------------------------------------
\subsection{7.1.4. チャネル}
\begin{frame}{7.1.4. チャネル: 画像は2次元ではなく3次元}
\begin{block}{これまでの単純化}
これまでは, 入力画像を2次元配列として
\[
  \X=\bigl([\X]_{i,j}\bigr)
\]
と書いていた.
しかし, 実際のカラー画像は赤, 緑, 青の3つのチャネルを持つ.
\end{block}

\begin{columns}[T,totalwidth=\textwidth]
\begin{column}{0.48\textwidth}
\begin{itemize}
  \item 画像の縦方向: 高さ.
  \item 画像の横方向: 幅.
  \item 各画素における成分: チャネル.
\end{itemize}

\[
  \Xten \in \R^{H\times W\times C}
\]
\[
  [\Xten]_{i,j,c}
\]
は, 位置 $(i,j)$ のチャネル $c$ の値.
\end{column}

\begin{column}{0.48\textwidth}
\centering
\begin{tikzpicture}[scale=0.62]
  % RGB sheets
  \foreach \s/\name/\col in {0/R/red!25,1/G/green!25,2/B/blue!25}{
    \begin{scope}[shift={(0.35*\s,0.25*\s)}]
      \draw[fill=\col,draw=gray!70] (0,0) rectangle (3.3,2.6);
      \draw[step=0.55,gray!55] (0,0) grid (3.3,2.6);
      \node at (3.75,2.15) {\name};
    \end{scope}
  }
  \draw[-{Latex[length=2mm]},thick] (4.4,1.5) -- (5.7,1.5);
  \node[align=center] at (7.4,1.5)
  {$\Xten=\bigl([\Xten]_{i,j,c}\bigr)$\\
   高さ $\times$ 幅 $\times$ チャネル};
\end{tikzpicture}
\end{column}
\end{columns}
\begin{block}{複数入力チャネルの場合}
入力が $[\Xten]_{i,j,c}$ のようにチャネル $c$ を持つため,
フィルタも $[\Vten]_{a,b,c}$ のようにチャネル方向の添字を持つ.
したがって, 空間方向 $(a,b)$ だけでなくチャネル方向 $c$ についても和を取る.
\end{block}
\end{frame}

% ------------------------------------------------------------
\begin{frame}{複数入力チャネルの畳み込み}
\begin{columns}[T,totalwidth=\textwidth]
\begin{column}{0.52\textwidth}
\begin{block}{1つの出力チャネルを作る場合}
入力チャネル数を $C_{\mathrm{in}}$ とすると,
\[
  [\Hid]_{i,j}
  =
  \sum_{a=-\Delta}^{\Delta}
  \sum_{b=-\Delta}^{\Delta}
  \sum_{c=1}^{C_{\mathrm{in}}}
  [\Vten]_{a,b,c}
  [\Xten]_{i+a,j+b,c}.
\]
\end{block}

\begin{itemize}
  \item $a,b$: 局所パッチ内の相対位置.
  \item $c$: 入力チャネル.
  \item $[\Vten]_{a,b,c}$: 位置とチャネルごとの重み.
\end{itemize}
\end{column}

\begin{column}{0.44\textwidth}
\centering
\begin{tikzpicture}[scale=0.62]
  % input stack
  \foreach \s/\col in {0/red!20,1/green!20,2/blue!20}{
    \begin{scope}[shift={(0.28*\s,0.18*\s)}]
      \draw[fill=\col,draw=gray!70] (0,0) rectangle (2.7,2.2);
      \draw[step=0.55,gray!55] (0,0) grid (2.7,2.2);
    \end{scope}
  }
  \draw[red,very thick] (0.55,0.55) rectangle (1.65,1.65);
  \node at (1.7,2.9) {入力 $\Xten$};

  % filter stack
  \foreach \s/\col in {0/red!35,1/green!35,2/blue!35}{
    \begin{scope}[shift={(4.2+0.28*\s,0.35+0.18*\s)}]
      \draw[fill=\col,draw=accentorange] (0,0) rectangle (1.1,1.1);
      \draw[step=0.55,accentorange!70] (0,0) grid (1.1,1.1);
    \end{scope}
  }
  \node at (5.0,2.4) {フィルタ $\Vten$};

  \draw[-{Latex[length=2mm]},thick] (6.3,1.2) -- (7.3,1.2);

  % output
  \draw[fill=softblue,draw=mainblue] (7.7,0.8) rectangle (8.3,1.4);
  \node at (8.0,2.1) {$[\Hid]_{i,j}$};
\end{tikzpicture}

\vspace{0.6em}
\small
RGBそれぞれの局所パッチに重みを掛けて, すべて足し合わせる.
\end{column}
\end{columns}
\begin{block}{隠れ表現にもチャネルを持たせる}
入力チャネルを $c$, 出力チャネルを $d$ とすると,
\[
  [\Hten]_{i,j,d}
  =
  \sum_{a=-\Delta}^{\Delta}
  \sum_{b=-\Delta}^{\Delta}
  \sum_c
  [\Vten]_{a,b,c,d}
  [\Xten]_{i+a,j+b,c}.
\]
\end{block}
\end{frame}

% ------------------------------------------------------------
\subsection{7.1.5. 要約と考察}

\begin{frame}{7.1.5. 要約: CNNは第一原理から理解できる}
\begin{enumerate}
  \item 画像を全結合層で扱うと，パラメータ数が大きくなりすぎる．
  \item 画像には，平行移動等価性と局所性という自然な構造がある．
  \item これらをMLPに制約として課すと，畳み込み層が現れる．
  \item 畳み込み層は，小さなカーネルを全位置で共有して特徴マップを作る．
  \item 複数チャネルにより，多種類の局所特徴を同時に学習できる．
\end{enumerate}

\vspace{3mm}
\begin{center}
\begin{tikzpicture}[>=Latex, node distance=8mm]
\node[draw, rounded corners, fill=softgray] (problem) {パラメータ爆発};
\node[draw, rounded corners, fill=softorange, right=of problem] (bias) {画像の帰納バイアス};
\node[draw, rounded corners, fill=softblue, right=of bias] (conv) {畳み込み層};
\draw[->, thick] (problem) -- (bias);
\draw[->, thick] (bias) -- (conv);
\end{tikzpicture}
\end{center}
\end{frame}

\section{7.2. 画像のための畳み込み}
\begin{frame}{7.2. 画像のための畳み込み}
\begin{block}{この節で学ぶこと}
前節では，画像に対して全結合層をそのまま使うのではなく，\alert{局所性}と\alert{重み共有}を使うと畳み込み層が自然に現れることを見た．
この節では，実際に畳み込み層がどのように計算されるかを，コードと数値例で確認する．
\end{block}

\begin{columns}[T,onlytextwidth]
\begin{column}{0.50\textwidth}
\begin{enumerate}
  \item 相互相関演算を手で計算する．
  \item それを \texttt{corr2d} として実装する．
  \item エッジ検出の例を見る．
  \item カーネルをデータから学習する．
  \item 特徴マップと受容野を理解する．
\end{enumerate}
\end{column}
\begin{column}{0.45\textwidth}
\centering
\begin{tikzpicture}[scale=0.64,>=Latex]
  \foreach \x in {0,...,5}{\foreach \y in {0,...,4}{\draw[mainblue!35] (\x,\y) rectangle +(0.75,0.75);}}
  \foreach \x in {1,2}{\foreach \y in {2,3}{\fill[accentorange!30] (\x,\y) rectangle +(0.75,0.75);\draw[accentorange,thick] (\x,\y) rectangle +(0.75,0.75);}}
  \node at (2.2,-0.55) {入力画像};
  \draw[->,thick] (5.0,2.2) -- (6.5,2.2);
  \foreach \x in {0,...,3}{\foreach \y in {0,...,2}{\draw[mainblue!35] (7+\x,1+\y) rectangle +(0.75,0.75);}}
  \fill[mainblue!40] (7.75,2.5) rectangle +(0.75,0.75);
  \node at (8.5,0.45) {特徴マップ};
\end{tikzpicture}
\end{column}
\end{columns}

\begin{seminaralertbox}
ここでの「畳み込み」は，深層学習の慣習では実際には\alert{相互相関演算}を指すことが多い．厳密な数学用語とは少しずれるので注意する．
\end{seminaralertbox}
\end{frame}

\subsection{7.2.1. 相互相関演算}

\begin{frame}{7.2.1. 相互相関演算：何を計算しているか}
\begin{block}{基本アイデア}
小さなカーネルを入力画像の上でスライドさせる．各位置で，\alert{重なった要素どうしを掛け算して足し合わせる}．その1つの値が出力の1画素になる．
\end{block}

\begin{center}
\begin{tikzpicture}[scale=0.72,>=Latex]
  % input 3x3
  \foreach \i/\row in {0/{0,1,2},1/{3,4,5},2/{6,7,8}}{}
  \foreach \x/\v in {0/0,1/1,2/2}{\node[draw,minimum size=8mm,fill=accentorange!20] at (\x,-0) {\v};}
  \foreach \x/\v in {0/3,1/4,2/5}{\node[draw,minimum size=8mm,fill=accentorange!20] at (\x,-1) {\v};}
  \foreach \x/\v in {0/6,1/7,2/8}{\node[draw,minimum size=8mm,fill=softgray] at (\x,-2) {\v};}
  \draw[accentorange,very thick] (-0.4,0.4) rectangle (1.4,-1.4);
  \node at (1,-2.8) {入力 $\X$};

  % kernel
  \node[draw,minimum size=8mm,fill=accentorange!20] at (4,0) {0};
  \node[draw,minimum size=8mm,fill=accentorange!20] at (5,0) {1};
  \node[draw,minimum size=8mm,fill=accentorange!20] at (4,-1) {2};
  \node[draw,minimum size=8mm,fill=accentorange!20] at (5,-1) {3};
  \node at (4.5,-2.0) {カーネル $\K$};

  \draw[->,thick] (2.0,-0.7) -- node[above]{要素ごとの積和} (3.3,-0.7);
  \draw[->,thick] (5.7,-0.7) -- (7.0,-0.7);

  % output
  \node[draw,minimum size=8mm,fill=mainblue!35] at (7.7,0) {19};
  \node[draw,minimum size=8mm,fill=softgray] at (8.7,0) {25};
  \node[draw,minimum size=8mm,fill=softgray] at (7.7,-1) {37};
  \node[draw,minimum size=8mm,fill=softgray] at (8.7,-1) {43};
  \node at (8.2,-2.0) {出力 $\Y$};
\end{tikzpicture}
\end{center}

\[
0\times0+1\times1+3\times2+4\times3=19.
\]
\end{frame}

\begin{frame}{出力サイズ：なぜ少し小さくなるのか}
\begin{block}{入力サイズとカーネルサイズ}
入力が $n_\mathrm{h}\times n_\mathrm{w}$，カーネルが $k_\mathrm{h}\times k_\mathrm{w}$ のとき，パディングなし・ストライド1の出力サイズは
\[
(n_\mathrm{h}-k_\mathrm{h}+1)\times(n_\mathrm{w}-k_\mathrm{w}+1).
\]
\end{block}

\begin{columns}[T,onlytextwidth]
\begin{column}{0.50\textwidth}
\begin{itemize}
  \item カーネル全体が画像の中に収まる位置でしか計算できない．
  \item そのため，画像の端では計算位置が減る．
  \item 例：$3\times3$ 入力に $2\times2$ カーネルを使うと，出力は $2\times2$ になる．
\end{itemize}
\end{column}
\begin{column}{0.45\textwidth}
\centering
\begin{tikzpicture}[scale=0.68]
  \foreach \x in {0,...,2}{\foreach \y in {0,...,2}{\draw[mainblue!45] (\x,\y) rectangle +(0.75,0.75);}}
  \draw[accentorange,thick] (0,1) rectangle +(1.5,1.5);
  \draw[accentorange,thick] (1.1,0.2) rectangle +(1.5,1.5);
  \node at (1.1,-0.55) {$3\times3$ 入力};
  \draw[->,thick] (3.2,1.1) -- (4.7,1.1);
  \foreach \x in {0,...,1}{\foreach \y in {0,...,1}{\draw[mainblue!45] (5+\x,0.5+\y) rectangle +(0.75,0.75);}}
  \node at (5.75,-0.05) {$2\times2$ 出力};
\end{tikzpicture}
\end{column}
\end{columns}

\begin{seminaralertbox}
出力サイズを小さくしたくない場合は，次節で学ぶ\alert{パディング}を使う．
\end{seminaralertbox}
\end{frame}


\begin{frame}[fragile]{\texttt{corr2d} の実装：出力サイズの準備}
\begin{lstlisting}
def corr2d(X, K):  #@save
    """Compute 2D cross-correlation."""
    h, w = K.shape  # kernel height and width

    # output size: (input size) - (kernel size) + 1
    Y = d2l.zeros((X.shape[0] - h + 1, X.shape[1] - w + 1))
\end{lstlisting}

\begin{block}{何をしているか}
入力 $X$ とカーネル $K$ を受け取り，
カーネルのサイズ $(h,w)$ を取り出す．

入力が $3\times 3$，カーネルが $2\times 2$ なら，
出力は
\[
(3-2+1)\times(3-2+1)=2\times2
\]
になる．
\end{block}
\end{frame}

\begin{frame}[fragile]{\texttt{corr2d} の実装：各位置で内積を計算}
\begin{lstlisting}
    # slide the kernel over all output positions
    for i in range(Y.shape[0]):
        for j in range(Y.shape[1]):

            # elementwise product, then sum
            Y[i, j] = d2l.reduce_sum((X[i: i + h, j: j + w] * K))

    return Y
\end{lstlisting}

\begin{block}{何をしているか}
\texttt{X[i: i+h, j: j+w]} は，入力 $X$ から
カーネルと同じ大きさの部分を切り出している．

その部分と $K$ を成分ごとに掛けて，すべて足すことで
\[
Y_{i,j}
=
\sum_{a=0}^{h-1}\sum_{b=0}^{w-1}
X_{i+a,j+b}K_{a,b}
\]
を計算している．
\end{block}
\end{frame}

\begin{frame}[fragile]{7.2.2 畳み込み層}
\begin{block}{畳み込み層の出力}
畳み込み層は，入力とカーネルの相互相関を計算し，さらにバイアスを加える．
\[
\text{出力}=\text{corr2d}(\text{入力},\ \text{カーネル})+\text{バイアス}.
\]
\end{block}

\begin{columns}[T,onlytextwidth]
\begin{column}{0.50\textwidth}
\begin{itemize}
  \item \alert{カーネル}：局所パッチを見るための重み．
  \item \alert{バイアス}：出力全体に足されるスカラー値．
  \item 学習では，カーネルとバイアスをデータから調整する．
\end{itemize}
\end{column}
\begin{column}{0.45\textwidth}
\centering
\begin{tikzpicture}[>=Latex,node distance=10mm]
\node[draw,rounded corners,fill=softgray,minimum width=24mm,minimum height=8mm] (x) {入力 $X$};
\node[draw,rounded corners,fill=softorange,below=of x,minimum width=24mm,minimum height=8mm] (k) {カーネル $K$};
\node[draw,rounded corners,fill=softblue,right=22mm of x,minimum width=28mm,minimum height=8mm] (y) {出力 $Y$};
\draw[->,thick] (x) -- node[above]{相互相関} (y);
\draw[->,thick] (k) -- (y);
\node[accentorange] at (3.6,-1.1) {$+b$};
\end{tikzpicture}
\end{column}
\end{columns}

\begin{seminaralertbox}
$h\times w$ 畳み込みとは，カーネルの高さが $h$，幅が $w$ である畳み込み層のこと．
\end{seminaralertbox}
\end{frame}

\begin{frame}[fragile]{畳み込み層の実装}
\begin{lstlisting}
class Conv2D(nn.Module):
    def __init__(self, kernel_size):
        super().__init__()

        # learnable kernel
        self.weight = nn.Parameter(
            torch.rand(kernel_size))

        # learnable bias
        self.bias = nn.Parameter(torch.zeros(1))

    def forward(self, x):
        return corr2d(x, self.weight) + self.bias
\end{lstlisting}
\begin{block}{ポイント}
\begin{itemize}
  \item \texttt{weight} が畳み込みカーネル．ここではランダムに初期化している．
  \item \texttt{bias} は最初は0．
  \item \texttt{forward} で \texttt{corr2d} を呼び出し，最後にバイアスを足す．
\end{itemize}
\end{block}
\end{frame}

\begin{frame}{7.2.3. エッジ検出：変化している場所を見つける}
\begin{block}{エッジとは}
画像の中で画素値が急に変わる場所を\alert{エッジ}と呼ぶ．
ここでは，白を $1$，黒を $0$ とし，中央が黒い簡単な画像を考える．
\end{block}

\begin{center}
\begin{tikzpicture}[scale=0.55]
\foreach \y in {0,...,5}{
  \foreach \x in {0,1,6,7}{\fill[gray!10] (\x,\y) rectangle +(0.8,0.8);\draw[gray!60] (\x,\y) rectangle +(0.8,0.8);}
  \foreach \x in {2,3,4,5}{\fill[gray!80] (\x,\y) rectangle +(0.8,0.8);\draw[gray!60] (\x,\y) rectangle +(0.8,0.8);}
}
\draw[accentorange,very thick] (1.6,0) rectangle (2.4,4.8);
\draw[accentorange,very thick] (4.8,0) rectangle (5.6,4.8);
\node at (3.2,-0.6) {白 $\to$ 黒，黒 $\to$ 白 の境界がエッジ};
\end{tikzpicture}
\end{center}
\end{frame}

\subsection{7.2.3. 画像における物体のエッジ検出}
\begin{frame}[fragile]{7.2.3 画像における物体のエッジ検出}
\begin{columns}[T,onlytextwidth]
\begin{column}{0.52\textwidth}
\begin{lstlisting}
X = d2l.ones((6, 8))
X[:, 2:6] = 0

K = d2l.tensor([[1.0, -1.0]])

Y = corr2d(X, K)
Y
\end{lstlisting}

\begin{seminaralertbox}
\[
K=\begin{pmatrix}1&-1\end{pmatrix}
\]
は横に隣り合う画素の差を取るカーネルである．
\end{seminaralertbox}
\end{column}

\begin{column}{0.45\textwidth}
\begin{block}{説明}
\begin{itemize}
\item $X$ は $6\times8$ の白黒画像である．
\item \texttt{X[:, 2:6] = 0} により，中央部分を黒くしている．
\item 左右で画素値が同じなら，差は $0$ になる．
\item $1\to0$ の境界では
\[
1\cdot1+0\cdot(-1)=1
\]
となる．
\item $0\to1$ の境界では
\[
0\cdot1+1\cdot(-1)=-1
\]
となる．
\item したがって，このカーネルは縦方向の境界を検出する．
\end{itemize}
\end{block}
\end{column}
\end{columns}
\end{frame}

\subsection{7.2.4 カーネルの学習}
\begin{frame}[fragile]{7.2.4 カーネルの学習}
\begin{columns}[T,onlytextwidth]
\begin{column}{0.56\textwidth}
\begin{lstlisting}
conv2d = nn.LazyConv2d(1, kernel_size=(1, 2),
                       bias=False)

X = X.reshape((1, 1, 6, 8))
Y = Y.reshape((1, 1, 6, 7))
lr = 3e-2

for i in range(10):
    Y_hat = conv2d(X)
    l = (Y_hat - Y) ** 2
    conv2d.zero_grad()
    l.sum().backward()
    conv2d.weight.data[:] -= lr * conv2d.weight.grad
\end{lstlisting}
\end{column}

\begin{column}{0.41\textwidth}
\begin{block}{説明}
\begin{itemize}
\item 前節では $K=(1,-1)$ を手で決めた．
\item ここでは，入力 $X$ と正解出力 $Y$ からカーネルを学習する．
\item \texttt{Y\_hat} は畳み込み層の予測出力．
\item \texttt{l = (Y\_hat - Y) ** 2} で二乗誤差を計算する．
\item 誤差が小さくなるようにカーネルを更新する．
\end{itemize}
\end{block}

\begin{seminaralertbox}
手で設計したエッジ検出カーネルも，
データから学習できる．
\end{seminaralertbox}
\end{column}
\end{columns}
\end{frame}

\begin{frame}[fragile]{学習されたカーネル}
\begin{columns}[T,onlytextwidth]
\begin{column}{0.50\textwidth}
\begin{lstlisting}
d2l.reshape(conv2d.weight.data, (1, 2))
\end{lstlisting}

\begin{block}{出力例}
\begin{lstlisting}
tensor([[ 0.9414, -1.0291]])
\end{lstlisting}
\end{block}
\end{column}

\begin{column}{0.47\textwidth}
\begin{block}{説明}
\begin{itemize}
\item 学習後の重みは
\[
(0.94,\,-1.03)
\]
のようになる．
\item これは，手で設定した
\[
K=(1,-1)
\]
に非常に近い．
\item つまり，ニューラルネットワークは
「エッジを検出するフィルタ」を自動で見つけたことになる．
\end{itemize}
\end{block}

\begin{seminaralertbox}
CNNでは，カーネルを人間が設計するのではなく，
学習によって獲得する．
\end{seminaralertbox}
\end{column}
\end{columns}
\end{frame}

\section{7.3. パディングとストライド}
\begin{frame}[fragile]{7.3 パディングとストライド}
\begin{columns}[T,onlytextwidth]
\begin{column}{0.50\textwidth}
\begin{block}{パディング}
\begin{itemize}
\item 畳み込みを行うと，出力サイズは小さくなる．
\item 境界の周りに $0$ を追加すると，出力サイズを保ちやすい．
\item これを\textbf{ゼロパディング}という．
\end{itemize}
\[
(n_h-k_h+p_h+1)
\times
(n_w-k_w+p_w+1)
\]
\end{block}

\begin{lstlisting}
conv2d = nn.LazyConv2d(1, kernel_size=3,
                       padding=1)
comp_conv2d(conv2d, X).shape
# torch.Size([8, 8])
\end{lstlisting}
\end{column}

\begin{column}{0.47\textwidth}
\begin{block}{ストライド}
\begin{itemize}
\item カーネルを何マスずつ動かすかを\textbf{ストライド}という．
\item 通常は $1$ マスずつ動かす．
\item ストライドを大きくすると，出力サイズを小さくできる．
\end{itemize}
\[
\left\lfloor
\frac{n_h-k_h+p_h+s_h}{s_h}
\right\rfloor
\times
\left\lfloor
\frac{n_w-k_w+p_w+s_w}{s_w}
\right\rfloor
\]
\end{block}

\begin{lstlisting}
conv2d = nn.LazyConv2d(1, kernel_size=3,
                       padding=1, stride=2)
comp_conv2d(conv2d, X).shape
# torch.Size([4, 4])
\end{lstlisting}
\end{column}
\end{columns}
\end{frame}

\begin{frame}{7.2--7.3 まとめ}
\begin{block}{7.2 畳み込み層}
\begin{itemize}
\item 畳み込み層では，入力画像 $X$ にカーネル $K$ を重ねて相互相関演算を行う．
\item 実際には
\[
Y = X \star K + b
\]
のように，相互相関の結果にバイアスを加える．
\item カーネル $K$ は，エッジなどの局所的な特徴を取り出すフィルタとして働く．
\item CNNでは，このカーネルを手で決めるのではなく，データから学習する．
\end{itemize}
\end{block}

\begin{block}{7.3 パディングとストライド}
\begin{itemize}
\item \textbf{パディング}は，入力の周囲に $0$ を追加して，出力サイズの縮小を抑える操作である．
\item \textbf{ストライド}は，カーネルを何マスずつ動かすかを決める量である．
\item ストライドを大きくすると，出力サイズは小さくなる．
\end{itemize}
\end{block}

\begin{seminaralertbox}
畳み込み層は「局所的な特徴を取り出す層」であり，
パディングとストライドによって出力サイズを調整できる．
\end{seminaralertbox}
\end{frame}




\end{document}
